\section{Why training fails}
\label{sec:mechanism}

We used two kinds of evidence: probes of nine saved checkpoints from the
controlled layer search (five failed, four working, plus two
single-condition references, where only $L$ or the seed changes), and
per-epoch records from six instrumented training runs
(Section~\ref{sec:dynamics}). We wanted to know what the internal features
of a failed model contain, what happens to its attention, and what separates
a failed seed from a working one at the same settings. Probe definitions,
per-checkpoint values and plots are in Supplementary Sections~S3--S4.

\subsection{What the internal features encode}
\label{sec:h1}

For each checkpoint we take the penultimate representation
$h\in\mathbb{R}^{256}$ on the test set, recover operating-condition labels by
$k$-means ($k{=}6$) on the three setting channels, and fit 5-fold
cross-validated linear and ridge probes for condition and RUL respectively.
Which condition an engine is running in is two to three times easier to read
off failed models (0.62--0.82) than off working ones (0.23--0.36), while RUL
readability falls from 0.66--0.76 to 0.04--0.09 on FD004 and from 0.78--0.81
to 0.55--0.58 on FD002 (Table~\ref{tab:mechanism}; full values in
Supplementary Table~S3 and t-SNE maps in Fig.~S2). The fixed output fits
this: per-condition label means on FD002 (86.7--87.1) and FD004 (92.7--93.3)
barely differ from the overall means, so features that carry only condition
have nothing better to output than the mean. On FD004 a failed model removes
the RUL signal from $h$ altogether; on FD002 a linear RUL component survives
even though the prediction is still 82.57, because the attention stack has
flattened the sequence. What fails is the usable representation, not simply
the information a linear probe can extract.

\subsection{What happens to attention}
\label{sec:h2}

For every layer of every checkpoint we compute the normalised entropy of the
causal attention distributions, averaged over heads, queries and test
samples, where 1.0 means every time step weighted equally. Every failed
checkpoint sits at 1.0000 to four decimals, with the largest weight at its
uniform value, while working models lie between 0.70 and 0.93
(Table~\ref{tab:mechanism}; Supplementary Table~S4 and Fig.~S3). Working
multi-condition models sharpen with depth (FD002-L3: 0.84, 0.77, 0.70),
concentrating on positions that matter for wear; failed ones never
distinguish positions. Flat attention on its own is not the problem: the
single-condition references (FD001 0.89, FD003 0.93) are the flattest
working models in the archive and predict well, because their TCN/GRU
residual branches carry the degradation signal regardless of attention. The
damaging combination is flat attention, features organised by condition and
a fixed output. Flat attention describes how the condition cue stops the
model from using time, rather than causing the failure by itself.

\begin{table*}[t]
\centering\footnotesize
\caption{Evidence from the probes, pooled over checkpoints: how well condition
and RUL can be read from the internal features, mean normalised attention
entropy, and output std. Per-checkpoint values are in Supplementary
Tables~S3--S4.}
\label{tab:mechanism}
\begin{tabularx}{\linewidth}{Xcccc}
\toprule
Checkpoints & Cond.\ acc. & RUL $R^2$ (FD002 / FD004) & Entropy & Pred.\ std \\
\midrule
Failed runs (5) & 0.62--0.82 & 0.55--0.58 / 0.04--0.09 & 1.0000 & $<10^{-3}$ \\
Working, multi-cond. (6) & 0.23--0.36 & 0.76--0.81 / 0.66--0.76 & 0.70--0.93 & --- \\
Working, single-cond. (2) & n/a & 0.81--0.87 & 0.89--0.93 & --- \\
\bottomrule
\end{tabularx}
\end{table*}

\subsection{What decides whether a run recovers}
\label{sec:dynamics}

The two sections above describe the end state. To see what separates a failed
seed from a working one at the same settings, six instrumented re-runs
recorded, each epoch, the gradient $L_2$ norms of the TCN, GRU, attention and
head groups, attention entropy, condition- and RUL-probe accuracy on a fixed
batch, and validation RMSE (protocol in Supplementary Section~S1). The
result changed our reading. A failed run has not fallen into a bad solution
of its own; it has failed to get out of a stage that every run passes
through. All six, including the five that ended well, first settle at the
constant output of Section~\ref{sec:collapse_def}: validation RMSE pinned at
$\operatorname{std}(y_{\mathrm{val}})$, entropy exactly 1.0, prediction std
near zero, for 16--34 epochs (Table~\ref{tab:escape}). What differs is what
happens next.

In all five working runs the recovery follows the same order. The
pre-clipping TCN gradient norm jumps by two orders of magnitude
($\times$129--235) within 1--3 epochs while the encoder gradient stays
flat; attention entropy drops below 1.0 for the first time (0.93--0.99) and
keeps falling; and condition-probe accuracy falls from 0.48--0.70 to
0.29--0.40 as the RUL probe rises. Before recovery the encoder gradient
decays monotonically ($102\to1.2$ over 37 epochs in FD002 $L{=}2$, seed 42),
since a fixed output gives attention no error signal. Recovery happens at
epochs 21--39 and overlaps the early-stopping horizon (patience 20 from a
flat stretch starting near epoch 5--10): working runs get there only because
tiny improvements in validation RMSE ($41.975\to41.971\to\dots$)
intermittently reset the patience counter. Archived failed runs stopped at
epochs 31--38, inside that window, so some may have been recoveries cut
short by the stopping rule. Two seeded re-runs show how much the seed
matters: FD002 $L{=}2$ recovered under seed 42 (19.04) but failed under seed
7 (43.87), and FD004 $L{=}1$, archived as failed, recovered under the
instrumented protocol (20.92), which also varied the train/val split. What
looked like settings that always fail are in fact settings where either
outcome is likely, depending on the seed. Because the flat stretch is
visible from epoch $\approx$5, tens of epochs before early stopping, it gives
a third check for the procedure of Section~\ref{sec:method_audit}.

\begin{table}[t]
\centering\footnotesize
\caption{When the six instrumented runs started to improve again. Flat and
Recovery are epoch numbers; TCN jump is the pre-clipping gradient-norm ratio
at that point.}
\label{tab:escape}
\begin{tabularx}{\linewidth}{XCCCCC}
\toprule
Run & Final RMSE & Outcome & Flat (ep.) & Recovery & TCN jump \\
\midrule
FD002 $L{=}2$, s42 & 19.04 & recovered & 34 & 39 & $\times$163 \\
FD002 $L{=}3$, s42 & 18.58 & recovered & 30 & 34 & $\times$150 \\
FD002 $L{=}2$, s7 & 43.87 & failed & 25 & never & --- \\
FD004 $L{=}1$, s42 & 20.92 & recovered & 16 & 21 & $\times$129 \\
FD004 $L{=}3$, s42 & 19.60 & recovered & 26 & 31 & $\times$235 \\
FD004 $L{=}3$, s123 & 20.49 & recovered & 32 & 37 & $\times$169 \\
\bottomrule
\end{tabularx}
\end{table}

\begin{figure}[t]
\centering
\includegraphics[width=0.72\linewidth]{figures/fig5_escape_race.png}
\caption{What separates a working run from a failed one, across six
instrumented runs. (a) Validation RMSE: all runs sit at
$\operatorname{std}(y_{\mathrm{val}})$ for a long stretch, and working runs
start improving at staggered epochs. (b) Attention entropy: exactly 1.0
during the flat stretch, dropping only when training recovers. (c) TCN
gradient norm (log scale): a hundredfold jump. (d) Condition-probe accuracy:
high during the flat stretch, falling sharply once training recovers.}
\label{fig:escape}
\end{figure}

\subsection{Putting the pieces together}
\label{sec:synthesis}

The evidence fits one account. From near-random initialisation, the quickest
way to reduce MSE before the inputs have been decoded is to output the label
mean; every audited multi-condition run reaches this fixed output within a
few epochs and stays there for 16--34 epochs. With six operating conditions
the features around it are organised by which condition the engine is in
(Section~\ref{sec:h1}); with a fixed output the attention branch gets no
error signal and ends up weighting every time step equally
(Section~\ref{sec:h2}); training then resumes only after a sudden jump in
the gradients of a residual branch, usually the TCN, which lets attention
work again and replaces the condition cue with task information
(Section~\ref{sec:dynamics}). Whether that happens before early stopping
depends on window length, depth and seed, which gives the pattern we see:
settings that train reliably, settings that usually fail (FD002
$L\in\{2,4\}$), and settings where the seed decides (FD004 $w{=}60$,
$L{=}3$: six of eight failed).

Two more observations complete the picture. Single-condition data is
unaffected because the operating settings are constant channels removed by
the variance filter: there is no condition cue to pick up, and the TCN/GRU
branches suffice even when attention is nearly flat (FD003: 0.93). The
component removals supply the causal complement (full table in
Supplementary Table~S6): taking out the TCN front end makes FD002 fail
(19.73 $\to$ 43.84) but only degrades FD004 ($\to$ 24.21), whereas taking
out the Transformer encoder or the difference features makes FD004 fail
($\to$ 43.68 / 43.73) and leaves FD002 well. Which branch the model depends
on changes with window length and condition count, and keeping more than one
usable branch is what makes training reliable. Fig.~\ref{fig:phase} sketches
the pattern that Section~\ref{sec:method} turns into rules.

\begin{figure}[t]
\centering
\includegraphics[width=0.60\linewidth]{figures/fig6_phase_diagram.png}
\caption{Which settings train reliably. Horizontal axis: how much the six
operating conditions change the sensor readings. Vertical axis: how often a
run starts improving again before early stopping (points at measured rates
with pooled run counts). The three regions are settings that train reliably
for every seed, settings that depend on the seed, and settings that usually
fail.}
\label{fig:phase}
\end{figure}
